\section{Polyglot and Ambiguous-Carrier Handling}
\label{sec:polyglots}

\subsection{Classification is not acceptance}

A file-type label selects work under policy. It is not an authorization
decision. Name evidence, caller-supplied media type, and bounded byte
recognition can disagree.
A separate authorizing parser can then expose a different whole-object
interpretation. Polyglot research shows that several file interpretations can
be combined and that detector behavior depends on the chosen grammar and
boundary~\cite{koch2025}. Parser differentials have the same
general consequence even when the carrier was not intentionally constructed as
a polyglot~\cite{jana2012,you2025}.

Let $R_p(x,n,m)$ be the eligible route set defined in
\cref{eq:eligible-routes}. Selection is admissible only when this set is a
one-element set:
\begin{equation}
  \mathsf{Route}_p(u)=s
  \quad\Longleftrightarrow\quad
  R_p(x,n,m)=\{s\}.
  \label{eq:type-convergence}
\end{equation}
Compatibility follows the versioned alias relation, not a textual MIME prefix.
For example, \texttt{audio/wav} and \texttt{audio/x-wav} can select one WAV
profile. A media suffix paired with an executable grammar is not such an alias.

A complete parser accounts for the whole object presented to that parser, but
source and output evidence have different domains. A selected source parser
accounts for ranges retained by structural reconstruction. Semantic
reconstruction may discard source bytes after bounded decoding.
The authorizing parser accounts for the reconstructed candidate \(y\). Its
success does not prove that every byte of source \(x\) had one interpretation.
Context can expose disguise but cannot override incompatible bytes. A handler's
declared output media type is a claim to test, not proof of the candidate format.

\subsection{Distinct ambiguous-carrier constructions}

The four carrier constructions differ in byte topology and therefore in the
observation most likely to falsify acceptance. Their comparison is shown in
\cref{fig:carrier-control-map}. Every route still passes the full predicate in
\cref{eq:extended-acceptance}. The rightmost column names the decisive
construction-specific control.

\begin{figure*}[t]
\centering
\resizebox{0.96\textwidth}{!}{%
\begin{tikzpicture}[
  cell/.style={draw,rounded corners=2pt,align=center,minimum height=9mm,font=\small,inner sep=3pt},
  kind/.style={cell,fill=blue!7,minimum width=33mm},
  bytes/.style={cell,fill=gray!8,minimum width=53mm,font=\small\ttfamily},
  control/.style={cell,fill=green!8,minimum width=58mm},
  arrow/.style={-{Latex[length=2mm]},thick}
]
\node[kind] (k1) at (0,0) {prefix camouflage};
\node[bytes] (b1) at (5.0,0) {[foreign prefix][media body]};
\node[control] (c1) at (11.5,0) {source evidence convergence;\\complete parse of candidate};
\node[kind] (k2) at (0,-1.3) {trailing object};
\node[bytes] (b2) at (5.0,-1.3) {[media terminal][suffix $z$]};
\node[control] (c2) at (11.5,-1.3) {source boundary accounted for or\\discarded by semantic reconstruction};
\node[kind] (k3) at (0,-2.6) {dual-format polyglot};
\node[bytes] (b3) at (5.0,-2.6) {[one byte string]\\$G_1,G_2:\mathsf{complete}$};
\node[control] (c3) at (11.5,-2.6) {conflicting source routes block;\\candidate is authorized separately};
\node[kind] (k4) at (0,-3.9) {nested active object};
\node[bytes] (b4) at (5.0,-3.9) {[container]\\$\supset$ [foreign component]};
\node[control] (c4) at (11.5,-3.9) {source component closure or\\bounded semantic reconstruction};
\foreach \i in {1,2,3,4} {
  \draw[arrow] (k\i) -- (b\i);
  \draw[arrow] (b\i) -- (c\i);
}
\end{tikzpicture}%
}
\caption{Carrier construction and decisive acceptance control. Prefix
camouflage, a suffix, a dual-format interpretation, and a nested component are
distinct test strata rather than interchangeable uses of ``polyglot.''}
\label{fig:carrier-control-map}
\end{figure*}

Complete consumption directly addresses trailing material in the candidate.
For every nonempty suffix \(z\in\mathbb O^*\) that an authorizing parser leaves after the
original terminal point,
\begin{equation}
\begin{aligned}
  &\mathsf{Parse}_{p,v,k}(y\parallel z)=\rho,\\
  &\rho.t=|y|<|y\parallel z|\\
  &\qquad\Longrightarrow
  \neg\mathsf{CompleteOK}_{p}(\rho,k,y\parallel z).
\end{aligned}
  \label{eq:suffix-rejection}
\end{equation}
This implication states non-acceptance by the authorizing parser. It does not infer
language nonmembership from parser failure. If the parser instead accepts
\(y\parallel z\) as a new complete object, all profile and component predicates
are evaluated on that entire object. Concatenation alone cannot establish
rejection. Declared-size containers apply the same accounting rule to RIFF, ISO
BMFF, and EBML parent boundaries.

Component closure rejects a nested foreign object only when it occupies a
forbidden structural position. It does not identify hidden meaning inside
retained pixels, samples, or codec frames. Prefix camouflage and true polyglots
add a source-evidence obligation: detected incompatible routes block. Semantic
reconstruction may emit a new object after one selected decoder consumes the
source under stated bounds, but acceptance of that output does not establish
that the source lacked another unrecognized interpretation. Transformation-time
decoder compromise remains governed by \cref{sec:process-isolation}.

\subsection{Reconstruction strength under ambiguity}

Semantic reconstruction can map several source representations into one output
language. A static image carrier is accepted only when the selected decoder
produces a bounded raster and the encoder creates a fresh PNG, JPEG, or WebP
candidate that passes the authorizing parser and a separate output decoder. GIF follows
the same principle at frame level. This supports the claim that source prefixes,
trailing objects, excluded container components, and excluded metadata were not
copied into the output. It does not prove that decoded pixels contain no
concealed information or that the selected decoder is universally correct.

Structural reconstruction has a narrower effect. APNG can preserve compressed
frame payloads. MP3 and FLAC can preserve encoded audio frames. WAV can preserve
PCM samples. Ogg can preserve media packets. ISO BMFF can preserve samples in
\texttt{mdat}, and WebM can preserve block frame payloads. Foreign structure in
discarded metadata or beyond an established boundary can be removed or rejected.
An alternate interpretation inside retained payload cannot be claimed
neutralized. A policy aimed at codec exploitation or payload-level overlap
therefore requires semantic reconstruction.

Even when authorizing parsing succeeds, publication still requires
\(\ell\geq q_{p,s}\) as defined in \cref{eq:level-satisfaction}. An APNG or
WebM structural candidate can therefore be well formed and remain
non-deliverable when policy requires semantic reconstruction. Canonical form
cannot silently weaken the requested reconstruction strength.

\subsection{Threat terminology and evidentiary boundary}

Malware labels are used only to organize scenario provenance, following the
behavioral distinctions in \cref{sec:background}. Method verdicts remain tied to
measurable carrier properties such as type disagreement, unsupported grammar,
complete-boundary failure, component closure, and reconstruction-level
insufficiency. Privacy-bearing metadata belongs to a separate leakage stratum
and is not evidence of spyware \cite{nist80083}.

Steganographic content and arbitrary data encoded in legitimate pixel or sample
values remain outside the structural polyglot claim. Semantic re-encoding can
preserve such information because media fidelity is part of the reconstruction
objective. Detecting covert meaning requires a separate content-analysis model,
corpus, and error specification.

The controlled evaluation therefore treats prefix camouflage, trailing objects,
true dual-format polyglots, and nested active objects as separate mutation
families. Structural and semantic policies are evaluated independently, with
the triggering carrier property recorded for each decision. Until the deferred
campaigns run, the justified contribution is the explicit decision rule and the
format-specific mechanism by which ambiguity is removed, rejected, or left
outside the claim.
